\specialchapter{序言}

本书汇集了迄今为止已在我的 Eléments de Mathématique 中发表过的大部分历史注记，未作实质性的修改。只是行文已使之独立于这些注记原先所附属的各卷 Eléments；因此，原则上，凡具备扎实的、相当于大学本科水平的经典数学修养的读者，都可以阅读本书。

当然，构成本书的这些专题研究，无论如何不能自诩勾勒出——哪怕是以粗略的方式——一部直至今天为止完整而连贯的数学发展史。经典数学的一些完整分支，如微分几何、代数几何、变分法，只是一笔带过；另一些分支，如解析函数论、微分方程论或偏微分方程论，则几乎未被触及；而且越是接近现代，这些空缺就越众多、越重要。不言而喻，这绝非出于故意的省略；其原因仅仅在于 Eléments 中相应的分册尚未出版。

最后，读者在这些注记中几乎找不到有关所论数学家的文献性或轶事性的资料；对每一门理论，本书首先力求尽可能清楚地揭示：其主导思想是什么，这些思想又是如何发展、如何相互作用的。

方括号中的数字指书末的参考文献。

\tableofcontents
\mainmatter

\chapter{数学基础；逻辑；集合论。}

对通常所说的“数学基础”的研究，自 19 世纪初以来一直在不间断地进行，但若没有一个使逻辑系统化的平行努力——至少是使逻辑中支配数学命题之间联系的那些部分系统化——这项研究是不可能取得成果的。因此，不可能把集合论的历史和数学形式化的历史同“数理逻辑”的历史分开来讲述。然而，传统逻辑如同近代哲学家的逻辑一样，其原则上的适用范围远远超出数学。所以读者不要指望在下文中找到一部逻辑学史，哪怕是以极其简略的形式写成的；我们只限于尽可能追溯逻辑的演变，而且仅限于它触及数学演变的那些部分。正是由于这个原因，对于非经典逻辑（多值逻辑、模态逻辑），我们将只字不提；我们更无法处理那些论争的历史，这些论争从智者派一直延续到维也纳学派，从未停止使哲学家们发生分裂——分歧既在于把逻辑应用于现实世界的对象或人类思想的概念是否可能，也在于如何应用。

在希腊时代之前就存在一种相当发达的数学，这在今天已无任何疑问。不仅整数和量的度量这些（已经非常抽象的）概念，在我们所掌握的来自埃及或迦勒底的最古老的文献中已被普遍使用，而且巴比伦代数凭借其方法的优雅与可靠，也不应被看作一堆靠经验摸索解出的问题的简单汇集。再者，虽然在这些文本中找不到任何与该词形式意义相近的“证明”，但有理由相信，这类解法——其一般性是通过特殊的数值应用显现出来的——的发现，若无最低限度的逻辑联系是不可能的（这些联系也许并非完全自觉，而颇像现代代数学家在着手一项计算时、在“形式地写下”全部细节之前所依赖的那些联系）（{[}232{]}，第~203~页及以下）。

希腊人本质的独创性，恰恰在于一种自觉的努力：把数学证明排成这样一个序列，使得从一个环节过渡到下一个环节不留下任何怀疑的余地，并迫使人们普遍同意。

希腊数学家在研究过程中，正如现代数学家一样，使用的是“启发式”的而非令人信服的论证，这一点，只要看看阿基米德的《方法》一文 {[}153 c{]}便可以证明（如果还需要证明的话）；其中还值得注意的是，它提到了早期数学家“已发现但未证明”的结果。\(^{1}\) 不过，从我们所知的最早一批详细文本（其年代为公元前 5 世纪中叶）起，数学文本的理想“典范”便已完全确立。它将在伟大的经典作家欧几里得、阿基米德和阿波罗尼奥斯那里达到最高的表现；在这些作者那里，证明的概念与我们的证明概念毫无二致。

我们没有任何文本可以据以追溯这种“演绎方法”的最初步骤；当我们刚一觉察到它的存在时，它在我们看来已近乎完善。人们只能认为，它相当自然地契合于希腊思想所特有的、对世界“解释”的永不停歇的追求，而这种追求早在公元前 7 世纪的爱奥尼亚哲学家那里便清晰可辨；此外，传统一致地把这种方法的发展和完善归功于毕达哥拉斯学派，其年代介于公元前 6 世纪末与公元前 5 世纪中叶之间。

正是这种对其目标和方法都有充分自觉的“演绎”数学，成为后世哲学思想和数学思想集中之所在。我们将看到，一方面，以数学为蓝本的“形式”逻辑逐步建立起来，并最终以形式化语言的创立而告终；另一方面，主要从 19 世纪初开始，数学的基本概念受到越来越多的质疑，人们将为澄清其本性作出巨大努力，尤其是在集合论诞生之后。

\section*{逻辑的形式化。}
\phantomsection
\addcontentsline{toc}{section}{逻辑的形式化。}

从我们掌握的有关公元前 5 世纪希腊哲学思想的（非常零散的）文本来看，总的印象是：这一时期被一种日益自觉的努力所支配，即把当时修辞学和数学如此成功地付诸实施的引导讨论的程序，推广到人类思想的全部领域——换句话说，去创立最一般形式的逻辑。哲学著作的语气此时发生了突变：在公元前 7 世纪和公元前 6 世纪，哲学家们只是断言或预言（至多依据同样含糊的类比勾画出含糊的论证），而从巴门尼德开始，尤其是到了芝诺，他们展开论辩，并力图提炼出可以作为其辩证法基础的一般原理：正是在巴门尼德那里可以找到排中律的最早陈述，而埃利亚的芝诺的那些“归谬法”证明则久负盛名。但芝诺写作于公元前 5 世纪中叶；而且，无论我们的文献有多少不确定之处，\(^{2}\) 很可能在这个时代，数学家们在他们自己的领域里已经在日常使用这些原理了。

如前文所说，追述这种逻辑在孕育过程中步步皆是的大量困难，以及由此产生的从埃利亚学派经智者派到柏拉图和亚里士多德的种种论争，并非我们的任务；我们在这里只拣出演讲艺术的勤勉耕耘及其后果之一——语言分析——在这一演变中所起的作用，这些发展通常被认为主要应归功于公元前 5 世纪的智者派。另一方面，数学的影响即使不总是被明确承认，也依然是显而易见的，在柏拉图和亚里士多德的著作中尤其如此。可以说，柏拉图几乎被数学迷住了；虽然他本人在这个领域并非创造者，但在他一生中的某个特定时期之后，他一直关注同时代数学家（其中许多人是他的朋友或学生）的发现，从此再未停止过一种最为直接的兴趣，甚至到了提出新研究方向的地步；因此，在他的著作中，数学不断地充当例证或模型（有时甚至像在毕达哥拉斯学派那里一样，滋养了他走向神秘主义的倾向）。至于他的学生亚里士多德，他不可能没有接受学园的学生所必需的最低限度的数学基础，人们甚至出版过一部由他著作中论及数学或涉及数学的文句辑成的集子 {[}153 d{]}；但他似乎从未下过大力气去了解他那个时代的数学动态，他在这个领域所引述的，只是很早以前就已普及的结果。这种错位在后世大多数哲学家那里只会愈加明显，他们当中许多人由于缺乏技术上的准备，真心诚意地自以为在内行地谈论数学，而实际上他们所指涉的不过是数学演进中早已被超越的一个阶段。

就逻辑而言，这一时期的最终产物是亚里士多德的宏篇巨制 {[}6{]}，其巨大功绩在于第一次成功地系统化并整理了在其前辈那里一直处于含糊或未定形的推理程序。\(^{3}\) 在这里，我们必须首先记住这部著作的总论点，即：一切正确的推理都可以化归为对一小批固定不变的规则的系统应用，而这些规则独立于所论对象的特殊本性（用字母来表示概念或命题——这大概是亚里士多德从数学家那里借用来的——清楚地表明了这种独立性）。但亚里士多德的注意力几乎完全集中在一种特殊类型的关系和逻辑连锁上，它们构成他所称的“三段论”：这主要是指一些若用集合论语言表述、今天会写成 \(A \subset B\) 或 \(A \cap B \neq \emptyset\) 形式的关系，\(^{4}\) 以及把这些关系或其否定借助如下格式串联起来的方式

\[
(A \subset B \text {   and   } B \subset C) \Rightarrow (A \subset C).
\]

亚里士多德对他那个时代的数学仍了解得足够多，因而看得出这类格式并不足以涵盖数学家所使用的全部逻辑运算，更不足以涵盖逻辑的其他应用（{[}6{]}，An. Pr., I, 46）。\(^{5}\) 至少，他对各种形式的“三段论”所作的深入研究（这一研究几乎全部用于澄清由推理所及对象的含糊或晦涩不断引起的种种困难）也使他有机会表述出求得一个命题的否定的规则（{[}6{]}，An. Pr., I, 46）。同样应归功于亚里士多德的，是他极精确地区分了“全称”命题与“特称”命题各自的作用，这是量词的最初雏形。\(^{6}\) 但是，他的著作（常被狭隘而缺乏理解地加以解释）直到 19 世纪很久以后仍有巨大影响，这种影响如何助长了哲学家对数学研究的忽视，并阻碍了形式逻辑的进步，这已是尽人皆知的事了。\(^{7}\)

然而，形式逻辑在古代仍继续取得进步，那是在作为逍遥学派对手的麦加拉学派和斯多葛学派的周围。遗憾的是，我们关于这些学说的信息全部来自二手，而且往往出自论敌或平庸注释者之口。这些逻辑学家取得的主要进展，看来在于建立了今天所说的“命题演算”：他们不像亚里士多德那样局限于 \(A \subset B\) 这种特殊形式的命题，而是对完全不确定的命题陈述规则。此外，他们对这些规则之间的逻辑联系分析得如此深入，以致知道如何从被立为“不可证明”的五条规则出发，用与现代方法非常类似的程序推出其余全部规则 {[}23{]}。可惜他们的影响相当短暂，他们的成果遂归于湮没，直到 19 世纪的逻辑学家重新发现它们为止。直到 17 世纪，亚里士多德始终是逻辑学界无可争议的宗师；众所周知，经院哲学家完全处于他的支配之下，而尽管他们对形式逻辑的贡献绝非微不足道，{[}25{]} 其中却不包含相对于古代哲学家成就而言在最高层次上的任何进步。

不过，在这里指出下面一点是恰当的：亚里士多德及其后继者的著作似乎并未对数学产生显著影响。希腊数学家沿着毕达哥拉斯学派及公元前 4 世纪后继者（泰奥多鲁斯、泰阿泰德、欧多克索斯）开辟的道路从事研究，而在表述其结果时显然并不拘泥于形式逻辑：只要把当时数学推理已经获得的灵活性和精确性，同亚里士多德逻辑尚处于非常初等的状况作一比较，这一事实就几乎不会令人惊讶。而当逻辑超越这一阶段时，指导其演变的又将是来自数学的新收获。

随着代数的发展，人们确实不可能不被形式逻辑规则与代数规则之间的类似所打动：二者共同的性质在于，都适用于（命题或数这样一些）并未精确确定的对象。而当 17 世纪代数记号在韦达和笛卡尔手中取得定型时，几乎立刻就出现了各种旨在表示逻辑运算的符号记号的尝试；但在莱布尼茨之前，这些尝试性的努力，例如埃里戈纳（1644 年）想用来记写初等几何证明的记号，或佩尔（1659 年）想用来记写算术证明的记号，都仍停留在非常表面的层次，没有带来数学推理分析上的任何进步。

到了莱布尼茨，我们面对的是一位同时也是第一流数学家的哲学家，他知道如何从自己的数学经验中提炼出使形式逻辑走出经院死胡同的思想萌芽。\(^{8}\) 作为一位举世罕见的博学通才、一个取之不竭的独创而多产的思想之源，莱布尼茨对逻辑的兴趣尤为浓厚，因为逻辑正处于他为使语言和思想形式化而设的宏伟计划的中心，而他毕生都在为这一计划不懈地工作。他自幼受经院逻辑的熏陶，却被这样一个（可上溯到雷蒙·吕勒的）想法所吸引：这种方法把一切人类概念还原为原始概念，构成一个“人类思想的字母表”，并以近乎机械的方式重新组合它们，从而得到一切真命题（{[}198 b{]}，第 VII 卷，第~185~页；参见 {[}69a{]}，第 II 章）。还在非常年轻的时候，他就形成了另一个独创得多的思想，即符号记号可以作为思想的“阿里阿德涅之线”而大有用处：\(^{9}\) “真正的 方法”，他说，“应当给我们一条阿里阿德涅之线，也就是说，一种粗略而有形的引导心灵的手段，就像几何中所画的线以及规定给算术学徒的那类运算一样。没有它，我们的心灵就无法沿着长路前行而不迷失。”（{[}198 b{]}，第 VII 卷，第~22~页；参见 {[}69 a{]}，第~90~页）。直到 25 岁上下，他对当时的数学都所知甚少，因此起初是以“通用语言”的形式提出他的计划（{[}69 a{]}，第 III 章）；但一与代数接触，他便把它采作其“普遍文字”的模型。他指的是这样一类符号语言：它能毫无歧义地表达一切人类思想，增强我们的推理能力，通过一种纯然机械的专注努力来避免错误，并且构造得使“连提出空想的人自己都不了解的那种空想，也不可能用这些字符写下来”（{[}198 a{]}，第 I 卷，第~187~页）。在莱布尼茨著作中论及这一宏伟计划及其实现将带来何种进步的无数段落里（参见 {[}69 a{]}，第 IV 章和第 VI 章），可以看到他对形式化语言的概念构想得何等清晰：一种纯粹的符号组合，其中重要的只是符号的连接，\(^{10}\) 以至一台机器就能产生出所有的定理，\(^{11}\) 而一切争论都可以通过简单的计算来解决（{[}198 b{]}，第 VII 卷，第~198-203~页）。即使这些希望今天看来有些过高，莱布尼茨数学工作的相当一部分——从他在无穷小演算的符号体系方面的工作开始（见第~190~页以下）——确实应当系于他思想中这种一以贯之的倾向，这一点仍然是真的；他自己对此有充分的自觉，还明确地把关于指数记号和行列式的想法（{[}198 a{]}，第 II 卷，第~204~页；参见 {[}69 a{]}，第~481-487~页）以及“几何演算”的概要（见第~50~页和第~61~页以下；参见 {[}69 a{]}，第 IX 章）也同他的“文字”联系在一起。但在他心中，本质的部分必须是符号逻辑，或者如他所说，是一个“Calculus ratiocinator”（推理演算），虽然他没有能创立这种演算，但我们至少看到他尝试了不下三次。第一次尝试时，他想到给每个“原始”词项关联一个素数，由若干原始词项组成的词项则用相应素数的乘积来表示；\(^{12}\) 他试图把通常的三段论规则翻译到这个系统中去，但遇到了由否定引起的重大复杂化（他很自然地想用变号来表示否定），于是很快放弃了这条道路（{[}198 c{]}，第~42-96~页；参见 {[}69 a{]}，第~326-344~页）。在后来的尝试中，他力图赋予亚里士多德逻辑一种更为代数化的形式；他有时保留记号 \(AB\) 来表示两个概念的合取，有时采用记号 \(A + B\)；\(^{13}\) 他（用乘法记号）写下了幂等律 \(AA = A\)，指出可以把命题“所有 A 都是 B”换成等式 \(A = AB\)，并且从这一点出发，用纯代数的计算就能重新得到亚里士多德的大部分规则（{[}198 c{]}，第~229-237~页和第~356-399~页；参见 {[}69 a{]}，第~345-364~页）；他还有了空集（“non Ens”）的想法，并且例如认识到了命题“所有 A 都是 B”与“A.(非 B) 不存在”的等价（出处同上）。此外，他注意到他的逻辑演算不仅适用于概念的逻辑，而且适用于命题的逻辑（{[}198 c{]}，第~377~页）。这样看来，他已经非常接近“布尔演算”。可惜的是，他似乎未能完全摆脱经院的影响；他不仅把自己演算的几乎唯一目标定为用他的记号去转录三段论的规则，\(^{14}\) 甚至为了完全恢复亚里士多德的规则而不惜牺牲自己最幸运的想法，哪怕这些规则中有些与空集观念不相容。\(^{15}\)

莱布尼茨的著作大部分直到 20 世纪初仍未发表，直接的影响因而甚微。整个 18 世纪以及 19 世纪初，有好几位作者（德·塞格纳、J. 兰伯特、普卢克凯、霍兰德、德·卡斯蒂隆、热尔戈纳）做了与莱布尼茨类似的尝试，但从未实质性地超出他停步的地方；这些工作收效甚微，以致他们中的大多数人对自己先行者的成果一无所知。\(^{16}\) 无论如何，G. 布尔——现代符号逻辑真正的创立者非他莫属 {[}29{]}——正是在同样的条件下从事写作的。他的关键想法在于系统地采取考虑“外延”的立场，即直接用集合来计算，把两个集合的交记作 xy，把 x 与 y 没有公共元素时二者的并记作 \(x + y\)。他还引入了一个记作 1 的“全集”（全体元素的集合）和记作 0 的空集，并把 x 的余集写作 1 - x。像莱布尼茨所做过的那样，他用关系 xy = x 来解释包含关系（他由此毫不费力地提取出对经典三段论规则的论证），而他对并与余集的记法也给他的系统带来了先行者们所缺乏的灵活性。\(^{17}\) 此外，通过把每个命题同它成立的那些“情形”的集合相关联，他把蕴涵关系解释为一种包含关系，于是他的集合演算便以这种方式为他提供了“命题演算”的那些规则。

19 世纪下半叶，布尔的系统成了一个活跃的逻辑学家学派的工作基础，他们在好几点上改进并完善了它。例如，杰文斯（1864 年）扩大了并运算 \(x + y\) 的含义，把它推广到 \(x\) 与 \(y\) 任意的情形；A. 德·摩根在 1858 年、C. S. 皮尔士在 1867 年证明了对偶律

\[
(\mathbf {C} A) \cap (\mathbf {C} B) = \mathbf {C} (A \cup B), (\mathbf {C} A) \cup (\mathbf {C} B) = \mathbf {C} (A \cap B); ^{18}
\]

德·摩根还在 1860 年着手研究关系，定义了二元关系的逆与复合（也就是与图上的运算 \(\overline{G}^{-1}\) 和 \(G_{1} \circ G_{2}\) 相对应的那些运算）。\(^{19}\) 所有这些工作都在施罗德卷帙浩繁而多产的巨著 {[}277{]} 中得到系统的阐述和发展。但颇为有趣的是，上面谈到的这些逻辑学家似乎对把自己的成果应用于数学并不很感兴趣；相反，布尔尤其是施罗德似乎以模仿经典代数的方法和问题（往往相当人为）来发展“布尔”代数作为自己的主要目标。这种态度的原因，无疑应当从如下事实中去寻找：布尔演算当时还缺乏转写大多数数学推理的手段，\(^{20}\) 因而对于莱布尼茨的伟大梦想只给出了一个很不完全的回答。建造一种更适合于数学的形式体系——其中变量和量词的引入（各自独立地归功于弗雷格 {[}117 a, b, c{]} 和 C. S. 皮尔士 {[}248 b{]}）构成主要的阶段——是逻辑学家和数学家的工作，与上述诸人不同，他们的首要目标恰恰是数学基础方面的应用。

弗雷格的计划 {[}117 b 与 c{]} 是要为算术建立一个基础，它基于一种借助“概念文字”（Begriffschrift）而形式化的逻辑；关于他定义自然数的方式，我们将在后面（第~29~页）再回来讨论。他的工作的特点，是对概念的分析极其精确、对细节极为讲究；正是由于这种倾向，他引入了许多后来在现代逻辑中显得极为重要的区分：例如，正是他第一个区分了命题的陈述与断言该命题为真，区分了属于关系与包含关系，区分了对象 x 与只含这个对象的集合 \(\{x\}\)，等等。他的形式化逻辑不仅包含数学中习用意义上的“变量”，还包含代表不确定关系的、可施以量词的“命题变元”，它后来（通过罗素和怀特海的工作）为元数学提供了基本工具。可惜他所采用的符号缺乏启发性，排版复杂得吓人，又远离数学家的实践；这就使数学家们望而却步，并大大削弱了弗雷格对其同时代人的影响。

皮亚诺的目标同时要宏大得多，也脚踏实地得多：出版一部完全用形式语言写成的“数学公式集”，其中不仅包含数理逻辑，而且包含数学各最重要分支的全部成果。他在一大群热心合作者（瓦伊拉蒂、皮耶里、帕多亚、瓦卡、维万蒂、法诺、布拉利-福尔蒂）的帮助下，以那样的速度完成了这一雄心勃勃的计划，本身就见证了他所采用符号体系的卓越：他的语言紧贴数学家的现行实践，引入了大量精心挑选的缩写符号，又特别由于一套用句点分隔符代替括号的巧妙办法，仍然相当易于阅读 {[}246 f{]}。皮亚诺创制的许多记号今天已被大多数数学家采用：我们举出 \(\in, \supset\)（但与当今的用法相反，其含义是“包含于”或“蕴涵”\(^{21}\)）、\(\cup, \cap, A - B\)（由差 \(a - b\) 组成的集合，其中 \(a \in A\)，\(b \in B\)）。另一方面，正是在《数学公式集》中，第一次出现了对函数一般概念以及直接像\(^{22}\) 与逆像概念的透彻分析，还有“序列不过是定义在 N 上的函数”这一评注。但是，量词在皮亚诺那里受到掣肘的限制（在他的系统中，原则上只能对形如 \(A \to B\)、\(A \Leftrightarrow B\) 或 \(A = B\) 的关系加量词）。此外，他的一些门徒近乎狂热的热情使他们大遭讥笑；尤其是 H. 庞加莱的批评（常常并不公正）给皮亚诺学派以沉重的打击，成为其学说在数学界传播的障碍。

有了弗雷格和皮亚诺，今天所使用的形式语言的基本要素便已齐备。流传最广的无疑是在罗素和怀特海的巨著《Principia Mathematica》中锤炼出来的那一种，它幸运地把弗雷格的精确与皮亚诺的便利结合在一起 {[}266{]}。现今大多数形式语言与它的区别，仅在于一些旨在简化其用法的次要改动。其中最巧妙的，我们举出关系的“函数式”书写（例如以 \(\in xy\) 代替 \(x \in y\)，这是武卡谢维奇想出来的，靠它括号可以完全省略）；但最有趣的无疑是希尔伯特对符号 \(\tau\) 的引入，它使人可以把量词 \(\exists\) 和 \(\forall\) 看作缩写记号，避免引入皮亚诺和罗素的“泛函”符号 \(\iota\)（它只适用于函数关系），最后还使人免去在集合论中陈述选择公理的必要（{[}163 a{]}，第 III 卷，第~183~页）。

\section*{数学中的真理概念。}
\phantomsection
\addcontentsline{toc}{section}{数学中的真理概念。}

数学家们一直确信，他们证明的是“真理”或“真命题”；这样的信念显然只能是感情上的或形而上学的，而且，站到数学的地盘上既不能为它辩护，也不能赋予它一个不沦为同义反复的意义。因此，数学中真理概念的历史属于哲学史而不属于数学史；但这一概念的演变对数学的演变有过无可否认的影响，因此我们不能对它保持沉默。

首先应当指出，要看到一位具有深厚哲学素养的数学家，与要看到一位通晓数学的哲学家同样罕见；数学家关于哲学问题的意见，即使这些问题涉及他们自己的领域，也多半是从可疑的来源二手、三手接受来的意见。但也恰恰因为如此，使数学史家感兴趣的正是这些平庸的意见，其程度至少不亚于对下述思想家独创见解的兴趣：如笛卡尔或莱布尼茨（这两人都是第一流的数学家）、柏拉图（他至少还跟得上他那个时代的数学）、亚里士多德或康德（对这后两位就不能这样说了）。

数学真理的传统概念可以一直追溯到文艺复兴时期。在这种观念里，数学家研究的对象与自然科学研究的对象之间没有多大差别；二者都是可知的，人都已经通过直觉和理性把握了它们；没有理由怀疑直觉或理性，它们只有在使用得当时才会出错。“必须”，帕斯卡说，“有一种全然错乱的精神，才会在如此重大、几乎不可能失手的原则上推理错误”（{[}244{]}，第 XII 卷，第~9~页）。笛卡尔在自家炉火旁确信，“历来只有数学家能够找到一些证明，也就是说一些确实可靠的理由”（{[}85 a{]}，第 VI 卷，第~19~页），而且（如果相信他的自述的话）这发生在他建成如下形而上学之前很久；在那形而上学中，“我先前当作规则的那件事，也就是我们能够设想得非常清楚、非常分明的一切对象都是真的这件事”，他说，“之所以得到保证，只是因为上帝存在，且他是一个完满的存在”（{[}85 a{]}，第 VI 卷，第~38~页）。如果说莱布尼茨反驳笛卡尔，说看不出如何才能认定一个观念是“清楚而分明的”，\(^{23}\) 那么他也认为，只要表述被理解了，公理就是定义的明显而无可争议的推论。\(^{24}\) 同样不应忘记，在当时的语言里，数学包含许多我们今天已不再承认其为科学的学问，有时甚至一直延伸到工程师的技艺；而在它们所激发的信心之中，它们应用于“自然哲学”、“机械技艺”和航海所取得的惊人成功占了很大一部分。

在这种看待事物的方式下，公理与演绎规则一样，不再容许被讨论或质疑；至多可以让每个人按照自己的偏好，去选择“以古人的方式”推理，还是任凭自己的直觉发挥。出发点的选择同样是个趣味问题，于是出现了为数众多的《Elements》“改写本”，其中《Elements》坚实的逻辑框架变成了一幅离奇的讽刺画；人们给出无穷小演算和理性力学的一些概述，号称是从一些建立得惊人糟糕的基本原理推演出来的；而斯宾诺莎或许也是出于诚意，把他的《伦理学》说成是“more geometrico demonstrata”（以几何方式证明的）。如果说在 17 世纪很难找到两位在任何问题上都意见一致的数学家，如果说论战日日发生、无休无止而尖刻激烈，真理的概念却始终不曾受到质疑。“关于每个对象只有一个真理”，笛卡尔说，“谁找到了它，谁就知道关于它所能知道的一切”（{[}85 a{]}，第 VI 卷，第~21~页）。

虽然希腊鼎盛时期关于这些问题的数学文本没有一份留存下来，但希腊数学家在这个问题上的观点很可能要细致得多。演绎规则只是凭借经验才被提炼到足以令人完全信服的地步；在它们可以被认为凌驾于一切讨论之前，人们必须经历的主要是摸索和悖理。而且，如果以为就连帕斯卡判定为最明显的那些“公理”（据他妹妹所传的说法，他童年时凭着万无一失的本能自己发现了它们）都不曾成为长期讨论的对象，那就误解了希腊人的批判精神，误解了他们对论辩和巧辩的嗜好。在一个严格说来并非几何学的领域里，埃利亚学派的悖论为我们保存了这类论战的一些痕迹；阿基米德在指出（{[}5 b{]}，第 II 卷，第~265~页）他的前辈们在许多情形使用了我们通常冠以他名字的那条公理时，补充说，借助这条公理所证明的东西“与不用它所证明的东西同样地被人们接受”，而他的结果只要在同样的基础上被接受就够了。柏拉图按照他的形而上学观点，把数学呈现为通达“自在真理”的途径，把数学处理的对象说成在理念世界中具有实在的存在；但他还是在《理想国》的一段著名文字中精确地刻画了数学方法：“凡是研究几何与算术的人 \ldots{} 都假定偶数与奇数、三种角；他们把它们当作已知的东西：一旦作了假定，他们就认为自己不再需要对自己或对别人说明它们，{[}认为它{]}对每个人都清楚；从那里出发，他们循序渐进，以便经由一致的同意达到他们的研究所提示的目标”（{[}250{]}，第 VI 卷，510 c-e）。因此，构成一个证明的，首先是一个提供任意起点的出发点（尽管“对每个人都清楚”），他说，在此之外人们不再试图深挖；然后是一场沿着一串中间阶段循序进行的讨论；最后，在每一步，都由对话者的同意来为推理的正确性作保。还必须补充的是：公理一经陈述，原则上便不再容许任何对直觉的新诉求；普罗克洛斯引述格米努斯时提醒说，“我们从这门科学的先驱者们那里学到，当涉及必须成为我们几何学说之一部分的推理时，不要把仅只貌似可信的结论放在眼里”（{[}153 e{]}，第 I 卷，第~203~页）。

因此，正是经验与批判的熔炉，造就了数学推理规则的提炼；而如果确实像有人以相当可信的方式论证过的那样 \([317\ d]\)，欧几里得《Elements》第八卷为我们保存了阿尔希塔斯算术的一部分，那么在那里看到那种相当迂腐的推理的生硬便不足为怪了——在任何一个“发现”了或自以为发现了“严格性”的数学学派里，这种生硬都不会不出现。但是，这些推理规则一旦进入数学家的实践，直到不久以前似乎从未受到过怀疑：在亚里士多德和斯多葛学派那里，虽然部分规则是通过推理格式从其他规则推出来的，初始规则却总被认为是自明的。同样，在回溯到那些在他们看来为他们时代的科学提供了坚实基础（例如他们在传统归于希俄斯的希波克拉底（约公元前 450 年）的第一部《Elements》中所提出的那些）的“假设”、“公理”和“公设”之后，古典时期的希腊数学家似乎把全部精力都投入到新结果的发现，而不是对这些基础的批判——在那个时候，这种批判不可能不是徒劳的；而且，撇开一切形而上学的考虑不谈，上文所引柏拉图的文本所见证的，正是数学家们在其科学基础上的这种普遍一致。

另一方面，希腊数学家似乎并不认为有可能解释清楚作为他们出发点的那些“初始概念”，如直线、曲面、量的比；他们如果给出其“定义”，那显然出于良心的不安，而对其效力不抱任何幻想。不言而喻，反过来，对于“初始概念”之外的其他定义（常被称为“名义定义”），希腊数学家和哲学家有着完全清楚的观念。正是在这种背景下，“存在”问题在数学中大概第一次明确地出现了。亚里士多德没有忽略指出：一个定义并不蕴涵被定义对象的存在，在此之外还必须有一个公设或一个证明。他的这个看法无疑来自数学家的实践；无论如何，欧几里得都注意在把这些概念引入论证时，公设圆的存在，并证明等边三角形、平行线、正方形等的存在（{[}153 e{]}，第 I 卷）；这些证明都是“作图”；换句话说，他以公理为依据，展示出一些数学对象，并证明它们满足他需要加以辩护的那些定义。

于是我们看到，古典时代的希腊数学达到的是一种经验性的确定性（不管它出自某位哲学家的什么样的形而上学根据）；既然不可设想推理规则会受到质疑，希腊科学的成功、再加上人们关于批判性修订并不合时宜的感觉，就在公理本身所激发的信心中占了很大一部分；这种信心近乎上个世纪人们对理论物理学诸原理所抱的那种（同样是几无限制的）信心。学园的格言“nihil est in intellectu quod non prius fuerit in sensu”（理智中无一物不是先在感觉中的）所提示的，无论如何正是这一点；而笛卡尔反对的恰恰是这句格言，认为它没有为笛卡尔希望从理性的运用中提取的东西提供足够牢固的基础。

必须等到 19 世纪初，才能看到数学家们从笛卡尔式的傲慢（更不必说康德或黑格尔的傲慢；后者对他那个时代的科学发议论未免姗姗来迟，一如时尚\(^{25}\)）中走出来，恢复到一种像希腊人那样恰当平衡的立场。对经典概念的第一次打击，是高斯、罗巴切夫斯基和鲍耶在本世纪初建立的非欧双曲几何。我们不打算在这里详细追述这一发现的来龙去脉，它是证明平行公设的无数次徒劳尝试的结果（见 {[}105 a 和 b{]}）。这一发现在当时对数学原理的影响，也许不像有时所说的那么深刻。它只是迫使人们放弃上个世纪对欧氏几何“绝对真理”的奢望，尤其是迫使人们放弃定义蕴涵公理的莱布尼茨式观点；公理不再显得全然“自明”，而是显得像一些假设，有待确定它们是否适合于对感官世界的数学表示。高斯和罗巴切夫斯基认为，各种可能几何之间的争论可以由经验来裁决（{[}206{]}，第~76~页）。黎曼的观点也是如此，他那篇著名的就职演讲《论构成几何基础的假设》旨在为各种自然现象提供一个一般的数学框架：“尚待解决的”，他说，“是这些假设在多大程度上、在什么范围内被发现为经验所证实的问题”（{[}259 a{]}，第~284~页）。但这是一个显然与数学不再有任何关系的问题；而且上述作者中似乎没有人质疑：即使一种“几何”不对应于任何实验的实在，它的定理也仍然是“数学真理”。\(^{26}\)

尽管如此，即便如此，这种信服肯定不应再归之于对经典“几何直觉”的无限制信任；黎曼试图对“n 重延展的流形”（他工作的对象）所作的描述，只是在为引入“局部坐标”寻求辩护时才依赖“直观的”\(^{27}\) 考虑；从那时起，他看来自觉已经站在了坚实的地面上，即分析的地面。但分析最终毕竟建立在实数概念之上，而这个概念直到那时还保持着非常直观的性质；函数论的进展在这个方面导致了一些非常令人不安的结果：随着黎曼本人关于积分的研究，更随着波尔查诺和魏尔斯特拉斯构造的无切线曲线的例子，数学的整个病理学开始了。一个世纪以来，我们见到的这一物种的怪物如此之多，以致有些见怪不怪了，必须把最离奇的畸形特征堆积起来，才能再让我们吃惊。但这些现象在 19 世纪大多数数学家身上引起的反应，从厌恶到惊愕不一：“直觉怎么会把我们欺骗到这种地步？”H. 庞加莱自问道（{[}251 d{]}，第~19~页）；埃尔米特则（不无一丝幽默，而这句名言的评注者们并非全都察觉到了这幽默）宣称，他“怀着恐惧与厌恶，从这场没有导数的连续函数的令人悲叹的瘟疫面前转过身去”（{[}160{]}，第 II 卷，第~318~页）。最糟糕的是，这些如此违背常识的现象，再也不能归咎于阐释得不好的观念，像“不可分量”时代那样（见第~173~页），因为它们在波尔查诺、阿贝尔和柯西的改革之后依然存在，而那次改革以与比例理论同样严格的方式确立了极限概念的基础（见第~154~页）。因此，这必须完全记在我们几何直觉的粗糙与不完备的账上；而从那以后，它作为一种证明手段被理所当然地贬黜，也是合理的。

这一认识必定会反过来作用于经典数学，而首当其冲的是几何学。无论欧几里得的公理构造曾经享有何种敬意，人们早已注意到其中不止一处缺陷，而且在古代即已如此。受到批评和证明尝试最多的，是平行公设；但欧几里得的追随者和注释者也曾试图证明其他公设（尤其是直角相等的公设），或者承认某些定义——例如直线或平面的定义——的不充分。16 世纪时，《Elements》的一位刊行者克拉维乌斯注意到缺少一条保证第四比例项存在的公设；莱布尼茨则指出，欧几里得使用几何直觉而不明确提及，例如当他承认（《Elements》第 I 卷命题 1）两个各自通过对方圆心的圆有一个公共点时（{[}198 b{]}，第 VII 卷，第~166~页）。高斯（他自己并不拒绝使用这类拓扑考虑）提请注意一个点（或一条直线）位于另外两个“之间”这一观念在欧氏作图中所起的作用，而这个观念却没有得到定义（{[}124 a{]}，第 VIII 卷，第~222~页）。最后，位移的使用——尤其是在“三角形相等的情形”中——长期被认为自明，\(^{28}\) 不久便在 19 世纪的批评者看来同样依赖于未曾陈述的公理。于是，在 1860 年至 1885 年间，出现了对几何学开端的各种局部修正（亥姆霍兹、梅雷、乌埃尔），旨在弥补其中一些缺陷。但只是到了 M. 帕施 {[}245{]}，放弃一切对直觉的诉求才成为一个正式提出并得到完全严格贯彻的纲领。他事业的成功很快为他引来了众多效法者，他们主要在 1890 年至 1910 年间，给出了彼此颇为不同的欧氏几何公理表述。这些著作中最著名的是皮亚诺用其符号语言写成的那些 {[}246 d{]}，尤其是希尔伯特的《Grundlagen der Geometrie》{[}163 c{]}，该书出版于 1899 年，以其论述的明晰与深刻，立即理所当然地成为现代公理学的宪章，甚至到了使其先行者湮没无闻的地步。的确，希尔伯特在那里不仅给出了欧氏几何的一个完备的公理系统，还把这些公理划分为不同类型的不同组，并致力于确定每一组公理的确切影响范围：不仅孤立地展开每一组的逻辑推论，而且讨论省略或修改其中某些公理所得到的各种“几何”（在这些几何中，罗巴切夫斯基几何和黎曼几何只作为特例出现\(^{29}\)）；他就这样在一个直到当时都被认为最接近感官实在的领域里，清楚地展示了数学家在选择公设时所享有的自由。尽管这些性质奇特的“元几何”使不止一位哲学家陷入窘迫，《Grundlagen》的论点还是很快被数学家们几乎一致地接受；H. 庞加莱尽管很难被指摘偏爱形式主义，也在 1902 年承认，几何公理是一些约定，对它们而言，通常所理解的“真理”概念不再有意义（{[}251 c{]}，第~66-67~页）。“数学真理”于是仅仅存在于从公理任意设定的前提出发的逻辑演绎之中。如下文将见（第~35~页以下），进行这些演绎所用的推理规则，其本身的有效性不久也将受到质疑，从而导致数学基本观念的一次彻底重塑。

\section*{对象、模型与结构。}
\phantomsection
\addcontentsline{toc}{section}{对象、模型与结构。}

\begin{enumerate}
\def\labelenumi{\Alph{enumi})}
\tightlist
\item
  数学的对象与结构。——从古代直到 19 世纪，关于数学家的主要对象是什么，存在着一种共同的理解；它们恰恰就是柏拉图在前文所引段落（第~13~页）中提到的那些：数、量与图形。如果说起初还必须加上作为力学、天文学、光学和音乐之题材的那些对象和现象，那么在希腊人那里，这些“数学性”的学科也总是与算术和几何清楚地分开，而从文艺复兴开始，它们相当快地取得了独立科学的地位。
\end{enumerate}

无论数学对象的概念被这样那样的数学家或哲学家的哲学细微色彩如何渲染，至少在一点上是一致的：这些对象是给定的，我们无权指派给它们任意的性质，正如物理学家不能改变一种自然现象一样。说实在的，心理反应无疑构成了这些看法的一部分，这些反应不是我们要追踪的，但每个数学家在徒劳地力图抓住一个似乎不断溜走的证明而耗尽心力时，都体会过它们。从这里到把这种阻力同感官世界横亘在我们路上的障碍等同起来，只有一步之遥；即使在今天，不止一位自称不妥协的形式主义者，在其内心深处也会乐于赞同埃尔米特的这段自白：“我相信，数和解析函数并不是我们心智的任意产物；我认为它们存在于我们之外，带有与客观实在的事物同样的必然性品格；我们遇见它们或发现它们，研究它们，就像物理学家、化学家和动物学家那样”（{[}160{]}，第 II 卷，第~398~页）。

在数学的经典观念里，根本谈不上游离于对数与图形的研究；但是，这个每一位数学家都认为自己必须在口头上表示赞同的官方学说，随着新观念的积累，却渐渐变成一种无法忍受的负担。代数学家面对负数的窘迫，直到解析几何给它提供了一个方便的“解释”才告结束；然而迟至 18 世纪，达朗贝尔在《百科全书》中（{[}75 a{]}，“NÉGATIF”条）讨论这个问题时，在一大段相当含混的解释之后突然没了底气，只满足于得出这样的结论：“关于负量的代数运算规则为人们普遍假定，并被普遍视为正确，无论人们对这些量另有何种想法。”至于虚数，丑闻要大得多；因为，如果它们是“不可能的”根，而且（直到 1800 年前后）看不到任何“解释”它们的途径，那么人们怎么能不出矛盾地谈论这些无法定义的存在，尤其是为什么要引入它们呢？达朗贝尔在这里保持着谨慎的沉默，甚至不曾提出这些问题，无疑是因为他意识到，自己无法作出比一个世纪前 A. 吉拉尔的天真回答更好的回答（{[}129{]}，f.~22）：“也许有人会问：这些不可能的解有什么用？我回答：有三点用处：为了保证一般规则，为了表明没有其他的解，也为了它本身的用处。”

在分析方面，17 世纪的局面也不见得好。一个幸运的情况是，解析几何仿佛应约而至地出现，以几何图形的形状为 17 世纪的伟大创造——函数概念——提供了一个“表示”，从而（与费马、帕斯卡和巴罗一道）有力地促成了无穷小演算的诞生（参见第~193~页）。但另一方面，人们也知道无穷小和不可分量这些观念将引起何等哲学-数学的论争。如果说达朗贝尔在这里要幸运些——他承认无穷小演算的“形而上学”中除了极限概念别无他物（{[}75 a{]}，“DIFFÉRENTIEL”条和“LIMITE”条，以及 {[}75 b{]}）——那么他仍不能比他的同时代人更明白发散级数展开的真实含义，也不能解释这样一个悖论：用剥除了一切数值解释的表达式算到最后，得到的却是精确的结果。最后，即使在“几何确定性”的领域内，欧几里得的框架也爆裂了：当斯特林在 1717 年不惮于说某条曲线有一个“二重虚无穷远点”时（{[}299{]}，新版第~93~页），要把这样一个“对象”与通常理解的概念联系起来，他肯定会有困难；而在 19 世纪初创立射影几何、使这类思想获得相当发展的彭赛列（见第~131~页），也仍然只满足于援引一个纯然形而上学的“连续性原理”来为自己辩护。

可以设想，在这些条件下（而且恰恰在“绝对真理”被以最强音宣布出来的那个时刻），证明的概念在 18 世纪中似乎变得越来越模糊，因为像希腊人那样，把进行推理所依据的概念及其基本性质一劳永逸地确定下来，已经无从谈起。始于 19 世纪初的向严格性的回归使这种状况有所改善，却没有因此挡住新观念的洪流：于是我们看到，代数中出现了伽罗瓦的虚数（{[}123{]}，第~113-127~页）、库默尔的理想数 {[}188 b{]}，继之而起的是向量和四元数、n 维空间、多重向量和张量（见第~61~页以下），更不消说布尔代数。无疑，重大的进步——正是它使严格性的回归得以不失去先前时代的任何收获——来自用更经典的术语为这些新观念给出“模型”的可能性：理想数或伽罗瓦虚数借助同余理论得到解释（见第~82~页以下）；n 维几何（如果愿意这样看）只表现为陈述“含 n 个变量”的代数结果的一种纯粹语言；至于经典的虚数——它以平面上点所作的几何表示（见第~161~页以下）标志着代数这场繁盛的开端——不久便可以在几何“模型”与用同余表述的解释之间任择其一（参见第~82~页）。但数学家们终于开始痛感，他们是在逆着一个自然的斜坡而奋斗，他们的工作正被这个斜坡拖着走；而在数学中，讨论没有任何合适“解释”的对象，应当是容许的：“数学的本质”，布尔在 1854 年说，“不在于埋头处理数与量的观念”（{[}29{]}，第 II 卷，第~13~页）。\(^{30}\) 同样的关切促使格拉斯曼在他的《Ausdehnungslehre》

1844 年版中，以一种起初排除了数或几何对象概念的形式来呈现他的演算。\(^{31}\) 而稍后，黎曼在他的就职演讲中，一开始就注意不说“点”，而在描述“n 重延展的流形”时说“确定方式”（Bestimmungsweise），并强调在这种流形中，“度量关系”（Massverhältnisse）“只能对抽象的量进行研究，只能用公式来表示；在一定的条件下，却可以把它们分解为一些关系，其中每一个孤立地看都容许一种几何表示，从而能够以几何的形式表达计算的结果”（{[}259 a{]}，第~276~页）。

从这时起，公理方法的扩张已是既成事实。如果说在一段时间里，人们还认为（只要可能）用几何直觉来检验“抽象”结果是有益的，那么至少已经公认，“经典”对象不再是数学家可以合法研究的唯一对象。这是因为——恰恰由于多种“解释”或“模型”皆属可能——人们已经认识到，数学对象的“本性”终究是次要的；比如，一个结果无论作为“纯”几何的定理呈现，还是借助解析几何作为代数定理呈现，都无关紧要。换句话说，数学的本质——这个直到那时只能用“通则”或“形而上学”之类含糊名称来表达的难以捉摸的概念——显现为对象之间关系的研究，而这些对象只是（有意地）通过其若干性质被了解和描述，恰好就是被当作公理置于其理论根基的那些性质。布尔在 1847 年就已经清楚地看到了这一点，他写道，数学处理的是“就其自身加以考虑的运算，独立于它们可以应用于其上的各种对象”（{[}29{]}，第 I 卷，第~3~页）。汉克尔在 1867 年开创代数的公理化时，捍卫的是一种“纯粹智性的数学，一种纯粹形式的理论，其目的不在于量、或量的影像即数的结合，而在于思想对象（“Gedankendinge”），可以有实在的对象或关系与它们对应，尽管这样的对应并非必要”（{[}146{]}，第~10~页）。康托尔在 1883 年以宣称“数学在其发展中是完全自由的，它的概念只受相容性这一必然性的约束，并通过精确定义与先前引入的概念相衔接”，来回应这一“自由数学”的主张（{[}47{]}，第~182~页）。最后，欧氏几何的修订成功地传播并普及了这些观念。帕施本人虽然仍执着于几何对象的某种“实在性”，也承认几何学实际上独立于这些对象的意义，纯粹在于研究它们的关系（{[}245{]}，第~90~页）；希尔伯特把这一观念推向其逻辑结论，强调甚至一个数学理论基本概念的名称也可以任意选择，\(^{32}\) 而庞加莱则把这一点表述为：公理是“伪装的定义”，从而完全颠倒了经院的观点。

因此，人们会禁不住说，现代的“结构”概念在实质上已于 1900 年前后达到；事实上，它还需要再经三十年左右的学徒期，才能以其全部光辉出现。当同类结构的性质相当简单时，辨认它们无疑并不困难；例如就群结构而言，这一阶段早在 19 世纪中叶便已达到。但与此同时，我们看到汉克尔在搏斗——却没有完全成功——以提炼出域和扩张的一般观念，他只能以一条半形而上学的“永久性原理”的形式来表达它 {[}146{]}，而这条原理要到 40 年后才由施泰尼茨 {[}294 a{]} 以最终的形式表述出来。尤其是，在这件事上，人们颇难摆脱这样一种印象：数学对象是连着其结构一起“给定”给我们的；只有对泛函分析相当长期的使用，才能使现代数学家熟悉如下想法：例如，有理数上有多个“自然的”拓扑，数直线上有多种测度。伴随着这种分离，向结构的一般定义的过渡终于实现了。

\begin{enumerate}
\def\labelenumi{\Alph{enumi})}
\setcounter{enumi}{1}
\item
  模型与同构。——用一门数学理论对另一门数学理论给出“模型”或“解释”这一观念的介入，我们已经多次注意到。这并不是一个新观念，而且我们在这里看到的，无疑是一种反复重现的感受的表现：诸门“数学科学”深层的统一性。如果可以把早期毕达哥拉斯学派的传统格言“万物皆数”视为可信，那么它可以被视为把当时的几何与代数归约为算术的第一次尝试的痕迹。尽管不可公度数的发现似乎永远关闭了这条道路，但它在希腊数学中引起的反应却是第二次综合的尝试，这一次以几何为基础，并把得自巴比伦人的代数方程解法等一并吸收进来。\(^{33}\) 众所周知，这一观念一直延续到 R. 邦贝利和笛卡尔的根本改革为止，后者把量的所有度量都同化为长度的度量（换句话说，同化为实数；参见第~151~页）。但随着笛卡尔和费马创立解析几何，倾向再次逆转，几何与代数得到了紧密得多的融合，不过这一次受益的是代数。此外，在别处，笛卡尔走得更远，构想出“所有这些通常被称为数学的科学 \ldots{} 尽管它们的题材各不相同”的本质统一性，“它们仍不会不汇合”，他说，“因为它们所考虑的，不过是其中可以找到的各种关系或比例”（{[}85 a{]}，第 VI 卷，第~19-20~页）。\(^{34}\) 然而这个观点只不过倾向于使代数成为基础的科学；莱布尼茨激烈地抗议这一结论，而他本人如前所述也构想了一种“普遍的数学”，但其规模宏大得多，而且已经相当接近现代观念。他把笛卡尔所说的“一致”精确化，实际上第一次窥见了同构的一般概念（他称之为“相似”），以及把同构的关系或运算“视为同一”的可能性；他给出了加法和乘法的例子（{[}69 a{]}，第~301-303~页）。但这些大胆的见解在他同时代人中没有引起任何回声，必须等到 19 世纪中叶前后代数的扩张（见第~51~页以下），才能看到莱布尼茨之梦开始实现。我们已经强调过，正是在那个时刻，“模型”成倍增加，通过简单的语言转换从一门理论过渡到另一门理论成为常事；其中最引人注目的例子也许是射影几何中的对偶（见第~132~页），当时把互为“对偶”的定理面对面排成两栏刊印的常见做法，无疑对同构概念的充分实现起了很大作用。从更技术的观点看，阿贝尔群的同构群概念当然已为高斯所知，置换群的同构群概念已为伽罗瓦所知（参见第~51~页以下）；对无论什么样的群，这一概念大约在 19 世纪中叶便已一般地获得。\(^{35}\) 此后，对每一门新的公理化理论，定义一个同构概念都是顺理成章的发展；但只是随着现代的结构概念，人们才最终认识到，每一结构都自带一个同构概念，无需对每种类型的结构分别给出专门的定义。
\item
  经典数学的算术化。——对“模型”观念越来越广泛的使用，还将使 19 世纪实现毕达哥拉斯学派所梦想的数学统一。在本世纪初，整数与连续量显得一如既往地不可调和，一如古代那样；实数仍然与几何量的概念（至少与长度的概念）联系在一起，而且人们正是求助于后者来构造负数和虚数的“模型”。甚至连有理数，也传统地附着于把一个量分成相等几份的想法；只有整数始终独立在外，正如高斯在 1832 年所说的那样，是“我们心智的专属产物”，他把它们与空间观念相对立（{[}124 a{]}，第 VIII 卷，第~201~页）。使算术与分析相互接近的最初努力，起初是围绕（正和负的）有理数进行的，出自马丁·欧姆（1822 年）；这些努力在 1860 年前后被几位作者重新提起，著名的有格拉斯曼、汉克尔和魏尔斯特拉斯（在其未刊的讲课中）；看来正是最后这一位想出了通过考虑整数对的类，来得到正有理数或负整数的“模型”的想法。但最重要的一步仍待迈出，即在有理数理论中为无理数找到一个“模型”；到 1870 年前后，这已成为一个紧迫的问题，因为在分析中发现“病态”现象之后，有必要把几何直觉的一切痕迹从实数定义中那个含糊的“量”的观念里清除出去。众所周知，这个问题大约就在那时几乎同时由康托尔、戴德金、梅雷和魏尔斯特拉斯用相当不同的方法解决了（见第~155~页）。
\end{enumerate}

从这时起，整数成了全部经典数学的基础。此外，随着公理方法的扩张和把数学对象视为心智的自由创造，基于算术的“模型”获得了更加重要的地位。对康托尔所宣称的这种自由，事实上仍然存在一个限制，那就是曾经萦绕于希腊人心头的“存在”问题；而它在这里以直接得多的方式出现，恰恰因为对直观表示的一切诉求此时都已被抛弃。我们将在后面（第~36~页以下）看到，“存在”概念在 20 世纪最初几年将成为怎样一个哲学-数学大漩涡的中心。但在 19 世纪，还没有走到那一步；要证明一个具有给定性质的数学对象存在，只需像欧几里得那样“构造”一个具有给定性质的对象。这恰恰正是算术“模型”的用处所在：一旦实数被“解释”成整数的说法，复数和欧氏几何也借助解析几何被如此解释，本世纪以来引入的所有新代数对象也都如此；最后——一项引起巨大轰动的发现——贝尔特拉米和克莱因甚至得到了罗巴切夫斯基和黎曼非欧几何的欧氏“模型”（见第~134~页），从而把这些乍看之下激起如此不信任的理论“算术化”了（并由此得到了完全的辩护）。

\begin{enumerate}
\def\labelenumi{\Alph{enumi})}
\setcounter{enumi}{3}
\tightlist
\item
  算术的公理化。——正是顺着这一演变，随后又转向了算术本身的根基，而事实上在大约 1880 年就可以看到这一点。看来在 19 世纪之前，除直接诉诸直觉之外，从未有人尝试定义整数的加法和乘法；莱布尼茨是唯一一个忠于自己原则、明确提醒人们 \(2 + 2 = 4\) 这样“自明”的真理只要想一想所涉各数的定义、便同样可以证明的人（{[}198 b{]}，第 IV 卷，第~403~页；参见 {[}69 a{]}，第~203~页）；而且他根本不认为加法和乘法的交换性是内禀的。\(^{36}\) 但他没有把这个想法再往前推进，而到 19 世纪中叶，这个方向上仍然毫无进展：就连其讲课大大助长了“算术化”观点传播的魏尔斯特拉斯，似乎也没有感到需要对整数理论作逻辑上的澄清。这个方向上的最初几步看来应归功于格拉斯曼，他在 1861 年（{[}134{]}，第 II 卷 \(_{2}\)，第~295~页）给出了整数加法和乘法的定义，并且只使用运算 \(x \mapsto x + 1\) 和归纳原理，证明了它们的基本性质（交换律、结合律、分配律）。归纳原理在 17 世纪已由 B. 帕斯卡清楚地构想并首次使用（{[}244{]}，第 III 卷，第~456~页）\(^{37}\)——尽管在古代也能找到它或多或少自觉的应用——而且从 17 世纪下半叶起就被数学家日常使用。但直到 1888 年，戴德金（{[}79{]}，第 III 卷，第~359-361~页）才陈述了算术的一个完备的公理系统（该系统三年后由皮亚诺重述，且通常冠以他的名字 {[}246 c{]}），其中特别包含归纳原理的一个精确表述（格拉斯曼使用了它，却没有明确陈述）。
\end{enumerate}

有了这种公理化，人们似乎已经达到了数学的最终基础。事实上，正当算术公理被清楚表述之时，对许多数学家（首先就是戴德金和皮亚诺本人）来说，这些公理就已经被剥夺了这门原始科学的角色，让位于数学诸理论中来得最晚的一个——集合论；而将要围绕整数概念展开的种种论争，是不能同 1900 年至 1930 年间那场大的“基础危机”分割开来的。
